\documentclass[11pt]{article}
\usepackage[margin=1in]{geometry}
\usepackage{amsmath,amssymb,graphicx,booktabs}
\usepackage{hyperref}

\title{Decision Transformer for Sequential Portfolio Allocation}
\author{Research Project}
\date{August 2026}

\begin{document}
\maketitle

\begin{abstract}
We frame multi-asset portfolio allocation as an offline sequence-modeling problem using a Decision Transformer (DT). Historical trajectories synthesized from 18 behavior policies train a causal transformer to predict simplex portfolio weights conditioned on Return-to-Go (RTG). On Universe~A (14 assets, 2000--2025) we train DT, transformer behavior cloning (BC), and six RL baselines across five seeds, then backtest on 2020--2025. No learned method beats buy-and-hold on Sharpe after 10\,bp proportional costs. DT tracks BC almost exactly (Sharpe $0.97\pm0.12$ vs.\ $0.98\pm0.09$), and an RTG-quantile sweep is essentially flat, indicating unused return conditioning. Momentum attains the highest test CAGR (27.1\%) at the cost of turnover and drawdown; PPO collapses to equal weights.
\end{abstract}

\section{Introduction}
Online reinforcement learning in financial markets is impractical due to transaction costs, non-stationarity, and the inability to replay history. Offline RL and sequence modeling offer an alternative: learn from fixed historical trajectories without environment interaction. The Decision Transformer~\cite{chen2021dt} reformulates RL as autoregressive modeling over $(\hat{R}_t, s_t, a_t)$ tuples, enabling conditioning on desired future returns at inference. We ask whether that conditioning improves sequential long-only allocation relative to classical policies and standard offline RL.

This paper specifies the allocation task as a finite-horizon MDP (Section~\ref{sec:mdp}), reports how each baseline is implemented and how it performs on the 2020--2025 holdout (Section~\ref{sec:baselines}), then describes the DT pipeline and remaining results. Section~\ref{sec:challenges} records engineering and scientific obstacles encountered in the work.

\section{Problem Definition}
\label{sec:mdp}
We treat daily long-only rebalancing as a discrete-time, finite-horizon Markov Decision Process
\[
\mathcal{M} = \langle \mathcal{S}, \mathcal{A}, P, R, \gamma, \rho_0, T_{\mathrm{ep}} \rangle
\]
driven by an \emph{exogenous} price process. The agent cannot affect next-day market returns; it only chooses weights and pays turnover costs.

\subsection{Time, assets, and prices}
Let $t=0,\ldots,T-1$ index NYSE trading days and $i=1,\ldots,N$ index investable assets. Universe~A uses $N=14$:
S\&P~500 ETF, Fidelity Growth, Apple, Microsoft, Amazon, Nvidia, JPMorgan, Walmart, Fidelity Energy, gold, silver, copper, a synthetic 10-year Treasury, and cash.
Cash is a unit-value index of T-bill yields converted to daily simple returns $r^{\mathrm{cash}}_t = \max(y^{\mathrm{bill}}_t/100/252,\,0)$. The bond is a duration approximation $r^{\mathrm{bond}}_t \approx -D\,\Delta y^{10\mathrm{Y}}_t + y^{10\mathrm{Y}}_{t-1}/252$ with $D=8$, compounded into a price index. Relative (gross) price returns are
\begin{equation}
  y_t = \Bigl( \tfrac{p_{1,t}}{p_{1,t-1}},\ldots,\tfrac{p_{N,t}}{p_{N,t-1}} \Bigr)^\top \in \mathbb{R}^N_{>0},
\end{equation}
clipped below at $10^{-8}$. Non-investable series (VIX, WTI, USD, term spreads, international indices) enter the state only.

\subsection{Action space}
The action $a_t = w_t \in \mathcal{A}$ is a target allocation on the probability simplex
\begin{equation}
  \mathcal{A} = \Delta^{N-1}
  = \bigl\{ w \in \mathbb{R}^N : w_i \ge 0,\; \textstyle\sum_{i=1}^N w_i = 1 \bigr\}.
\end{equation}
Shorting and leverage are forbidden. Implemented actors emit softmax (or Dirichlet-mean) vectors; any residual mass is renormalized. Initial holdings at episode start are equal weight $w_0 = \mathbf{1}/N$.

\subsection{State space}
The implemented observation is \emph{not} $(w_{t-1}, X_t)$. Previous weights do not enter the feature vector. Let $Z_t \in \mathbb{R}^{d_z}$ concatenate, on day $t$:
per-asset one-day returns; MACD histogram, RSI(14), Bollinger \%b(20), realized vol at 20 and 60 days; cross-sectional momentum ranks at 60 and 120 days; macro (VIX level, USD return, 10Y--3M and 30Y--5Y spreads, WTI and natural-gas returns); and binary event-window flags. Missing entries are filled with 0. With lookback $L=20$,
\begin{equation}
  s_t = \mathrm{vec}\bigl(Z_{t-L+1:t}\bigr) \in \mathbb{R}^{d_s},\qquad d_s = 2580.
\end{equation}
Mean and standard deviation of $s_t$ are estimated on the training split only and applied to all dates.

\subsection{Transition}
Given a realized price path $\{y_t\}$, the transition is deterministic in the observation:
\begin{equation}
  P(s_{t+1}\mid s_t, a_t) = \mathbf{1}\bigl[s_{t+1}=s(t+1)\bigr].
\end{equation}
$a_t$ does not affect $s_{t+1}$. The environment stores the last target $w_t$ for the cost term only; it does \emph{not} drift holdings with $y_t$ before the next rebalance. Turnover is therefore $\lVert a_t - a_{t-1}\rVert_1$ against the previous \emph{target}, not against mark-to-market weights $w'_{t} \propto w_{t-1}\odot y_t$. That is a modeling choice (and a limitation).

The episode terminates at $t=T_{\mathrm{ep}}$ (training: $T_{\mathrm{ep}}=252$; evaluation: the length of the date split). Terminal transitions do not bootstrap ($d_t=1$).

\subsection{Reward}
Proportional transaction cost $\mu=0.001$ (10\,bp of $L_1$ turnover) yields the net logarithmic return
\begin{equation}
  R(s_t,a_t) = R_t
  = \ln\bigl(\max\{ a_t^\top y_t - \mu \lVert a_t - a_{t-1}\rVert_1,\; \varepsilon \}\bigr),
\end{equation}
with floor $\varepsilon=10^{-8}$ so that $\log$ is defined if costs exceed gross return. The objective is undiscounted finite-horizon return $\mathbb{E}[\sum_{t=0}^{T_{\mathrm{ep}}-1} R_t]$ ($\gamma=1$ for DT/RTG). Value-based RL agents internally bootstrap with $\gamma=0.99$.

\subsection{Return-to-Go}
For a trajectory of rewards $(R_0,\ldots,R_{T_{\mathrm{ep}}-1})$,
\begin{equation}
  \hat{R}_t = \sum_{t'=t}^{T_{\mathrm{ep}}-1} R_{t'}
  \qquad (\gamma=1).
\end{equation}
DT consumes the standardized scalar $(\hat{R}_t - \mu_{\mathrm{RTG}})/\sigma_{\mathrm{RTG}}$ with $\mu_{\mathrm{RTG}}=-0.002$, $\sigma_{\mathrm{RTG}}=0.139$ from the training set. At inference the budget starts at a chosen quantile of $\{\hat{R}_0\}$ and updates $\hat{R}_{t+1}=\hat{R}_t-R_t$, resetting every 252 steps on multi-year tests.

\subsection{Initial distribution and offline data}
$\rho_0$ is uniform over valid training start dates (episode length 252, stride determined by sampling 200 windows per behavior policy). The offline dataset $\mathcal{D}=\{\tau_j\}_{j=1}^{3600}$ contains format-v2 tuples $(s,a,R,\hat{R})$ sharing a global state matrix. Learning is purely offline: no further interaction with $P$ during training except that PPO/SAC/A2C take single-step updates on buffer transitions (not a live market loop).

\section{Data}
Daily closes from the 30-year market file (1995--2025) are aligned to days where S\&P~500 is non-null, forward-filled at most 5 sessions, then restricted to Universe~A from 2000-09-01 through 2025-12-30 (6{,}369 trading days). Splits: train 2000--2016, validation 2017--2019, test 2020--2025 (1{,}507 test days). Walk-forward uses six 2-year blocks; folds that start on or before 2016-12-31 overlap training and are flagged in-sample.

\section{Initial Baselines}
\label{sec:baselines}
All policies map $(t,s_t)$ to $a_t\in\Delta^{N-1}$ and are scored with the same reward, 2\% annualized risk-free rate for Sharpe/Sortino, and 252-day year. Classical methods use only $y_{<t}$ (no learned parameters). Learned methods see $s_t$ from Section~\ref{sec:mdp}.

\subsection{Classical implementations}
\paragraph{Buy-and-hold / equal-weight.}
Both emit $a_t=\mathbf{1}/N$ every day. Because turnover is measured against the previous \emph{target} (always $\mathbf{1}/N$), both have $\approx 0$ turnover and therefore identical backtests on this cost model---they are the same policy under our MDP. A mark-to-market buy-and-hold would drift and differ; we do not simulate that.

\paragraph{Momentum.}
Lookback 60 days, long the top $k=\lfloor N/3\rfloor=4$ assets by trailing cumulative gross return $\prod y-1$, equal-weight among winners, else $\mathbf{1}/N$ if $t<60$.

\paragraph{Mean-variance (max-Sharpe).}
On a 252-day window, estimate $\mu$ and Ledoit--Wolf covariance $\Sigma$, then 200 projected-gradient steps on the simplex: $w \leftarrow \Pi_{\Delta}(w - \eta(\mu-\Sigma w))$ with $\eta=0.05$. Fallback $\mathbf{1}/N$ until the window is full.

\paragraph{Min-variance.}
Same optimizer with gradient $2\Sigma w$ (variance, not Sharpe).

\paragraph{Risk parity.}
Inverse-volatility weights $w_i \propto 1/(\mathrm{std}_{60}(y_{\cdot,i})+10^{-8})$ on a 60-day window.

\paragraph{Behavior mixture (training data only).}
The offline buffer also includes Dirichlet-noisy copies of MVO, momentum, and equal-weight ($\alpha\in\{5,15,50\}$) and random Dirichlet($\alpha\in\{0.5,1,2\}$), 200 windows each, 18 policies, 3{,}600 episodes. These noisy policies are not evaluated as test strategies.

\subsection{Classical results (test 2020--2025)}
Table~\ref{tab:test} is the baseline report. Buy-and-hold/equal-weight Sharpe is $1.12$ (CAGR $23.7\%$, max DD $-23.3\%$). Mean- and min-variance are statistically indistinguishable from that (Sharpe differences $-0.003$, bootstrap CIs covering 0). Momentum has the best CAGR ($27.1\%$) with Sharpe $0.97$, DD $-32.0\%$, and $15.1\%$ mean daily turnover. Risk parity is a failed baseline (CAGR $0.6\%$, Sharpe $-0.18$): inverse-vol on this mix (cash, bond, commodities, megacap growth) over-weights the low-vol sleeve and does not capture the 2020--2025 equity run.

\begin{table}[h]
\centering
\caption{Test period (2020--2025), Universe~A, classical baselines. Turnover is mean daily $L_1$ turnover.}
\label{tab:test}
\begin{tabular}{lrrrr}
\toprule
Strategy & CAGR & Sharpe & Max DD & Turnover \\
\midrule
Momentum       & 0.271 & 0.973 & $-0.320$ & 0.151 \\
Mean-variance  & 0.242 & 1.116 & $-0.233$ & 0.007 \\
Buy-and-hold   & 0.237 & 1.118 & $-0.233$ & $\sim 0$ \\
Equal-weight   & 0.237 & 1.118 & $-0.233$ & $\sim 0$ \\
Min-variance   & 0.235 & 1.116 & $-0.233$ & $\sim 0$ \\
Risk parity    & 0.006 & $-0.178$ & $-0.233$ & 0.001 \\
\bottomrule
\end{tabular}
\end{table}

\subsection{Learned baseline implementations}
All use the same 5 seeds $\{42,123,456,789,1011\}$, 50 epochs, 200 steps/epoch, batch 256, and simplex outputs. Actors/critics are two-layer MLPs (width 256) unless noted. Checkpoints store every \texttt{nn.Module} under a \texttt{modules} dict.

\paragraph{Transformer BC.}
Same GPT as DT ($K=30$, $d_{\mathrm{model}}=192$, 4 layers, 6 heads) but tokens are $(s,a)$ only. MSE on the last-step action.

\paragraph{TD3+BC.}
Twin critics, delayed actor, BC MSE on the dataset action, Polyak $\tau=0.005$, $\gamma=0.99$. Target $r+\gamma(1-d)\min_i Q_{\theta'_i}(s', \pi(s'))$.

\paragraph{IQL.}
Expectile value $V$ fitted to a \emph{target} $Q$; critic bootstraps $V(s')$; advantage-weighted BC actor; Polyak target $Q$.

\paragraph{CQL.}
Twin $Q$ plus CQL(H) penalty: $\mathrm{logsumexp}$ over 10 random simplex actions per state (not over the batch). Actor: $-\,Q$ plus BC MSE.

\paragraph{PPO / SAC / A2C (reference).}
Single-step updates on offline transitions, not environment interaction. PPO and SAC use a Dirichlet actor (concentration $\mathrm{softplus}(f(s))+1$); inference uses the Dirichlet mean. A2C uses softmax. These are not a fair online comparison.

\subsection{Learned baseline results}
Table~\ref{tab:learned} reports mean $\pm$ std over seeds (DT included for comparison; it is not a classical baseline).

\begin{table}[h]
\centering
\caption{Test period (2020--2025), learned methods. Mean across five seeds; Sharpe std in parentheses. DT is conditioned at the 90th-percentile training RTG.}
\label{tab:learned}
\begin{tabular}{lrrrr}
\toprule
Strategy & CAGR & Sharpe & Max DD & Turnover \\
\midrule
PPO$^\dagger$ & 0.237 & $1.118$ $(0.000)$ & $-0.233$ & $\sim 0$ \\
SAC           & 0.222 & $1.039$ $(0.143)$ & $-0.248$ & 0.058 \\
BC            & 0.219 & $0.976$ $(0.095)$ & $-0.267$ & 0.059 \\
DT            & 0.215 & $0.966$ $(0.119)$ & $-0.272$ & 0.061 \\
IQL           & 0.196 & $0.818$ $(0.183)$ & $-0.303$ & 0.127 \\
CQL           & 0.343 & $0.781$ $(0.191)$ & $-0.499$ & 0.195 \\
TD3+BC        & 0.190 & $0.660$ $(0.234)$ & $-0.351$ & 0.135 \\
A2C           & $-0.014$ & $-0.127$ $(0.338)$ & $-0.539$ & 1.139 \\
\bottomrule
\multicolumn{5}{l}{\footnotesize $^\dagger$PPO is identical to buy-and-hold on every seed (Dirichlet mean collapsed to uniform).}
\end{tabular}
\end{table}

PPO's Dirichlet mean is $\mathbf{1}/N$ on every seed---it reproduces the strongest classical baseline by collapsing, not by allocating. SAC is the only learned agent with mean Sharpe above 1.0 ($1.04\pm0.14$) and still trails buy-and-hold. BC is a tight clone of the behavior mixture (Sharpe $0.98\pm0.09$) and is the right control for DT. IQL and TD3+BC lose on both return and risk. CQL's mean CAGR ($34.3\%$) is the headline high-water mark but with Sharpe $0.78$, vol $34.6\%$, and max DD $-50\%$ (seed 789 alone has $53\%$ CAGR). A2C is unusable (daily turnover $>1$).

\section{Decision Transformer}
A causal GPT interleaves $(\hat{R},s,a)$ tokens. The action head is softmax. Loss is MSE on the last context step. Training matches the learned baselines (50 epochs $\times$ 200 steps, batch 256, AdamW $10^{-4}$, gradient clip 1.0). Inference zero-fills the current-step action slot (the causal mask already hides it). Default target is the 90th percentile of training $\hat{R}_0$.

\section{Experiments}
\textbf{Metrics:} CAGR, Sharpe, Sortino, Calmar, maximum drawdown, mean daily turnover, VaR/CVaR.

\textbf{Inference:} Each of 40 checkpoints is rolled out on the test window. Walk-forward uses six 2-year folds.

\textbf{Statistics:} Circular block-bootstrap 95\% CIs (block size 20, 1{,}000 replicates) and Ledoit--Wolf tests versus buy-and-hold. Deflated Sharpe~\cite{bailey2014dsr} uses $n_{\mathrm{trials}}=46$ (every classical row and every seed).

\section{Additional Results}

\subsection{DT versus BC and buy-and-hold}
DT does not improve on BC (Sharpes $0.966$ vs.\ $0.976$). Neither beats buy-and-hold. Block-bootstrap CIs for DT versus buy-and-hold exclude zero (DT worse) on two of five seeds; the other three include zero. With $n_{\mathrm{trials}}=46$, every deflated Sharpe is $\approx 0$.

\subsection{RTG calibration}
Sweeping target RTG over training-episode quantiles $\{0.1,0.25,0.5,0.75,0.9\}$ (budget reset every 252 days) yields nearly identical realized log returns (mean $\approx 1.17$). Mean Sharpe by quantile is $1.01$, $1.00$, $1.00$, $0.97$, $0.97$---slightly inverse. Conditioning on a higher desired return does not produce a higher realized return.

\subsection{Walk-forward (out of sample)}
Table~\ref{tab:wf} shows Sharpe on folds that begin after the training cutoff. Learned policies follow buy-and-hold through 2022--2023 and 2024--2025 but never lead. Momentum wins 2020--2021 and loses 2022--2023.

\begin{table}[h]
\centering
\caption{Out-of-sample walk-forward Sharpe. Learned entries are means across five seeds.}
\label{tab:wf}
\begin{tabular}{lrrrrrr}
\toprule
Fold & BAH & Mom.\ & DT & BC & SAC & CQL \\
\midrule
2018--2019 & 0.80 & 0.86 & 0.79 & 0.78 & 0.79 & 0.19 \\
2020--2021 & 1.24 & 1.32 & 1.07 & 1.06 & 1.09 & 0.95 \\
2022--2023 & 0.33 & 0.07 & 0.25 & 0.28 & 0.33 & $-0.02$ \\
2024--2025 & 1.94 & 1.38 & 1.84 & 1.85 & 1.87 & 1.44 \\
\bottomrule
\end{tabular}
\end{table}

\section{Challenges}
\label{sec:challenges}

\paragraph{Exogenous, cost-aware MDP versus a strong $1/N$.}
Because $P$ does not depend on $a_t$ and $\mu>0$ taxes turnover, any deviation from a flat book must buy enough edge to cover costs. On 2020--2025 that bar was not cleared by DT, BC, or the offline RL suite. The coincidence of buy-and-hold and equal-weight under target-based turnover further inflates the baseline: a drifted buy-and-hold was never measured.

\paragraph{RTG unused.}
The scientific premise is return conditioning. Calibration is flat and DT $\approx$ BC, so the extra RTG channel is not used. Contributing factors include: (i)~standardizing $\hat{R}$ with a small $\sigma_{\mathrm{RTG}}=0.14$ while test paths last $\sim$1{,}500 days; (ii)~resetting the budget every 252 days, which severs a six-year target; (iii)~MSE-to-behavior-actions, which encourages cloning $\mathcal{D}$ rather than changing $a$ with $\hat{R}$.

\paragraph{Correctness of offline RL.}
Early agents used $s'=s$ (self-referential Bellman targets), never updated target networks, and computed CQL's $\mathrm{logsumexp}$ over the batch instead of over actions. IQL's $V$ chased live $Q$ (loss $10^4$). PPO Dirichlet concentrations hit NaN; A2C broke on post-v2 batch shapes. Checkpoints originally stored the algorithm \emph{name} rather than weights. These were fixed (real $s'$, Polyak, CQL(H), BC-anchored CQL actor) before the sweep; without them the learned table would be meaningless.

\paragraph{Systems constraints.}
A GTX~1070 (8\,GB, Pascal, no tensor cores) saturates at $\sim$1{,}300 samples/s. A full pass is $\sim$10.4~min/epoch; 5 seeds $\times$ 2 transformers $\times$ 50 full epochs is $\sim$87 hours, so training uses 200 steps/epoch ($\sim$3 data passes). GPU-resident buffers plus batch 256 caused OOM until storage defaulted to pinned CPU with optional residency. Right-aligned padding could mask an entire attention row (softmax NaN); the diagonal is now forced attendable. Index sampling originally drew a full \texttt{randperm} per batch ($O(\text{steps}\cdot n)$). Mixed precision was skipped: GP104 fp16 is far slower than fp32.

\paragraph{Evaluation gaps that shaped the run.}
Until checkpoint policies existed, tables contained only classical rows. Ablation runs overwrote \texttt{dt\_seed42.pt} across $K$. Universe~B shared \texttt{data/processed/}, so a B prepare would have destroyed A. The sweep is now resumable; this paper's numbers are Universe~A only (no $K$ or $\mu$ ablation, no B training).

\paragraph{Statistics.}
Counting every seed as a trial ($n=46$) makes deflated Sharpe vacuously zero, including for buy-and-hold. That is conservative to the point of being uninformative; an algorithm-level $n_{\mathrm{trials}}$ would be more appropriate and was not used here.

\paragraph{Data.}
Close-only series: no volume or microstructure. WTI on 2020-04-20 is negative and is a feature, not an investable asset. Vectorized Dirichlet sampling for noisy policies is statistically equivalent to per-day \texttt{dirichlet} but not bit-identical.

\section{Discussion}
On this close-only, long-only universe, a cost-aware $1/N$ policy is a formidable baseline. Sequence models reproduce it with extra turnover and do not extract a Sharpe premium from RTG. Offline RL either collapses (PPO), underperforms (TD3+BC, IQL), or inflates return by taking more risk (CQL). The ``online'' agents are buffer updates, not live interaction.

\textbf{Limitations} not already in Section~\ref{sec:challenges}: Universe~B (Bitcoin + aluminum from 2014-09-17) and $K,\mu$ sweeps were out of scope for this run.

\textbf{Future work.} Universe~B, context-length and cost ablations, drifted (mark-to-market) buy-and-hold, RTG losses that penalize ignoring $\hat{R}$, and algorithm-level deflated Sharpe.

\begin{thebibliography}{9}
\bibitem{chen2021dt}
L.~Chen, K.~Lu, A.~Rajeswaran, K.~Lee, A.~Grover, M.~Laskin, P.~Abbeel, A.~Srinivas, and I.~Mordatch, ``Decision Transformer: Reinforcement Learning via Sequence Modeling,'' NeurIPS, 2021.
\bibitem{bailey2014dsr}
D.~H.~Bailey and M.~L\'opez de Prado, ``The Deflated Sharpe Ratio: Correcting for Selection Bias, Backtest Overfitting and Non-Normality,'' \emph{J.~Portfolio Management}, 2014.
\end{thebibliography}

\end{document}
